\paragraph{Flujos de detritos y de barro.} Un aluvión es una mezcla de agua y sedimento cuyo comportamiento depende de la concentración volumétrica total de sólidos $\Cvt$. Con $\Cvt$ entre $0{,}20$ y $0{,}55$ la mezcla presenta un esfuerzo de fluencia $\tau_y$: se detiene cuando el esfuerzo motriz no lo supera. Su densidad es $\rho_m=\rho_w+\Cvt(\rho_s-\rho_w)$, donde $\rho_m$, $\rho_w$ y $\rho_s$ son las densidades de la mezcla, el agua y los sólidos, respectivamente (\si{kg.m^{-3}}); $\rho_m$ se sitúa entre \SI{1300}{kg.m^{-3}} y \SI{1900}{kg.m^{-3}}. En las relaciones reológicas interviene, en cambio, la concentración de la matriz fina $\Cvf$ (sección~\ref{sec:modelo}).

\paragraph{Rangos de concentración y densidad en la literatura.} Las crecidas con barro y los flujos hiperconcentrados tienen entre un \SI{20}{\percent} y un \SI{45}{\percent} de sedimento en volumen, con promedios de \SI{20}{\percent} a \SI{35}{\percent} y picos cercanos al \SI{45}{\percent} en los flujos de barro \citep{obrien1988,flo2d2021}; su densidad queda entre unos \SI{1330}{kg.m^{-3}} y \SI{1740}{kg.m^{-3}}. Sobre un \SI{50}{\percent} en volumen se habla de flujo de detritos \citep{costa1984,zanuttigh2007}, con concentraciones de sólidos de \SI{50}{\percent} a \SI{80}{\percent} y densidades de \SI{1400}{kg.m^{-3}} a \SI{2400}{kg.m^{-3}} medidas en canales experimentales y en muestras naturales \citep{iverson1997}; en oleadas densas observadas en terreno se han estimado unos \SI{2000}{kg.m^{-3}} \citep{zanuttigh2007}. La tabla~\ref{tab:densidades} sitúa en ese marco los barros ensayados en este estudio.

\begin{table}[htbp]\centering\small
\caption{Concentración volumétrica de sedimento ($C_v$) y densidad de la mezcla según la literatura, y barros ensayados en este estudio. Las densidades de los barros ensayados se calculan con $\rho_s=\SI{2650}{kg.m^{-3}}$; $\tau_y$ y $\eta$ son la fluencia y la viscosidad de la matriz fina según \citet{obrien1988}.}\label{tab:densidades}
\begin{adjustbox}{max width=\textwidth}\begin{tabular}{lll}\toprule
Tipo de flujo o barro & $C_v$ & $\rho_m$ (\si{kg.m^{-3}})\\\midrule
Crecida con barro y flujo hiperconcentrado \citep{obrien1988,flo2d2021} & 0,20--0,45 & \numrange{1330}{1740}\\
Flujo de barro, valores típicos \citep{flo2d2021} & 0,20--0,35 (picos de 0,45) & \numrange{1330}{1740}\\
Umbral de flujo de detritos \citep{costa1984,zanuttigh2007} & $\geq0{,}50$ & $\geq\num{1825}$\\
Flujo de detritos \citep{iverson1997} & 0,50--0,80 & \numrange{1400}{2400}\\\midrule
Este estudio: Glenwood~2b ($\tau_y=\SI{6,6}{Pa}$, $\eta=\SI{0,77}{Pa.s}$) & 0,40 & \num{1660}\\
Este estudio: Glenwood~2b con $C_v=0{,}43$ ($\tau_y\approx\SI{11}{Pa}$) & 0,43 & \num{1710}\\
Este estudio: Glenwood~1 ($\tau_y=\SI{29}{Pa}$, $\eta=\SI{8,8}{Pa.s}$) & 0,45 & \num{1742}\\
\bottomrule\end{tabular}\end{adjustbox}
\end{table}

\paragraph{Elección del barro y su motivo.} Tres evidencias orientan la elección. Según BioBioChile, Sernageomin clasificó el evento como una crecida repentina que arrastró gran cantidad de sedimento \citep{biobio2026origen}; en las imágenes del día consultadas se interpreta una matriz fina con piedras, troncos y vehículos arrastrados, sin frentes de bloques visibles; y el alcance observado, hasta la Av.\ Padre Hurtado, exige un barro capaz de recorrer varios kilómetros por calles de baja pendiente. Para el modelo se adopta un rango entre el extremo alto de los flujos de barro y el umbral de los flujos de detritos. El barro de referencia (Glenwood~2b, $C_v=0{,}40$) está en ese rango; en las simulaciones, uno más espeso (Glenwood~1, $C_v=0{,}45$) se detiene antes de Padre Hurtado, y el de referencia se extiende más allá de la zona afectada oficial. Como la fluencia crece exponencialmente con la concentración \citep{obrien1988}, un cambio pequeño ---de 0,40 a 0,43--- casi duplica $\tau_y$; por eso se ensayó ese valor intermedio, que además representa la resistencia adicional de la carga gruesa (piedras, troncos y escombros) que la reología de la matriz fina no incluye.

\paragraph{Aguas someras con reología.} Cuando el espesor del flujo es pequeño frente a su extensión, las ecuaciones de Navier–Stokes integradas en la profundidad se reducen a las de Saint-Venant en dos dimensiones. La reología de la mezcla entra como una pendiente de fricción $S_f$, definida en la ecuación~\eqref{eq:friccion}. El cierre cuadrático de \citet{obrien1988} y \citet{obrien1993}, usado por FLO-2D \citep{flo2d2021}, suma tres términos: fluencia, viscosidad y turbulencia más colisión de granos.

\paragraph{Representación urbana.} En ciudad, el flujo se canaliza por las calles y lo desvían los edificios. Estos pueden representarse como bloques elevados, como huecos impermeables o mediante porosidad; la elección condiciona la capacidad predictiva \citep{schubert2012}.

\paragraph{Volumen movilizable.} El volumen de un flujo de detritos puede estimarse con un balance entre lluvia y escorrentía. En cuencas recién quemadas existen además modelos empíricos que lo relacionan con la intensidad de lluvia de \SI{15}{min}, el área quemada y el relieve \citep{gartner2014}. El caudal pico se relaciona empíricamente con el volumen \citep{mizuyama1992}.

\paragraph{Calibración con observaciones cualitativas.} Cuando solo se conoce la extensión del área afectada, distintas combinaciones de volumen, caudal y resistencia producen huellas parecidas, es decir, hay equifinalidad. El enfoque de estimación de incertidumbre mediante verosimilitud generalizada (GLUE, por sus siglas en inglés; \citealp{beven1992}) propone aceptar todas las simulaciones de comportamiento aceptable e informar la envolvente de sus predicciones en lugar de un único mejor ajuste.
